\subsection{7. Was die vier Rechnungen gemeinsam haben und was
nicht}\label{was-die-vier-rechnungen-gemeinsam-haben-und-was-nicht}

{\def\LTcaptype{none} % do not increment counter
\begin{longtable}[]{@{}
  >{\raggedright\arraybackslash}p{(\linewidth - 6\tabcolsep) * \real{0.2500}}
  >{\raggedright\arraybackslash}p{(\linewidth - 6\tabcolsep) * \real{0.2500}}
  >{\raggedright\arraybackslash}p{(\linewidth - 6\tabcolsep) * \real{0.2500}}
  >{\raggedright\arraybackslash}p{(\linewidth - 6\tabcolsep) * \real{0.2500}}@{}}
\toprule\noalign{}
\begin{minipage}[b]{\linewidth}\raggedright
Rechnung
\end{minipage} & \begin{minipage}[b]{\linewidth}\raggedright
Tatsächlicher Träger
\end{minipage} & \begin{minipage}[b]{\linewidth}\raggedright
Neue Voraussetzung
\end{minipage} & \begin{minipage}[b]{\linewidth}\raggedright
Nicht übertragen
\end{minipage} \\
\midrule\noalign{}
\endhead
\bottomrule\noalign{}
\endlastfoot
Zwei Charts & 128 Fermionmoden, zwei Bosonbanken, N=2 & Überlappung c
und globale Vertexsumme & Keine native Auswahl von c \\
Ein normierter C-Tensor & 64L Fermionmoden, eine Bosonbank, N=2 und 3 &
Normiertes symmetrisches C & Keine unveränderte lokale
Einzelkopienamplitude \\
Phasen-/Recorderprozess & Originale einzelne Bank, N=2, K3 plus Zeiger &
Schaltbarkeit, Eingang, Detektor, Arbeit & Kein
N=64-Grundzustandsversuch \\
Lorentz-Auslese & Vorgeschriebene Impuls- und Feldfasern &
Lorentzadapter und Kinetikvertrag & Keine aus W hergeleitete Raumzeit \\
\end{longtable}
}

Die gemeinsame Quelle W und der gemeinsame Operator Nb erzeugen
nachvollziehbare Anschlüsse. Sie machen die vier Modelle nicht
nachträglich zu einem einzigen bereits ausgewählten Universum.
Insbesondere wurde der frühere native Grundzustand bei N=64 nicht für
die neuen globalen Kompositionen bewiesen. Seine v1.6.8-Polgrenzen
bleiben als historische, modellgebundene Ergebnisse erhalten, ohne neue
Verbesserung in dieser Runde.

\subsection{8. T1-T8: konkrete Bewegung der
Grenze}\label{t1-t8-konkrete-bewegung-der-grenze}

Die Benennungen folgen der Arbeitszuordnung der vollständigen v1.6.8.
\textbf{Kein vollständiges T-Gate wird geschlossen.}

{\def\LTcaptype{none} % do not increment counter
\begin{longtable}[]{@{}
  >{\raggedright\arraybackslash}p{(\linewidth - 4\tabcolsep) * \real{0.3333}}
  >{\raggedright\arraybackslash}p{(\linewidth - 4\tabcolsep) * \real{0.3333}}
  >{\raggedright\arraybackslash}p{(\linewidth - 4\tabcolsep) * \real{0.3333}}@{}}
\toprule\noalign{}
\begin{minipage}[b]{\linewidth}\raggedright
Front
\end{minipage} & \begin{minipage}[b]{\linewidth}\raggedright
Beitrag dieser Runde
\end{minipage} & \begin{minipage}[b]{\linewidth}\raggedright
Erstes weiter fehlendes Stück
\end{minipage} \\
\midrule\noalign{}
\endhead
\bottomrule\noalign{}
\endlastfoot
T1 Ursprung/Auswahl & Clock-Rang 60; zwei exakte
Kompositions-Gegenmodelle; inverse Antworten & Typisierte ursprüngliche
globale Einbettung und ausführbare Kontrollen \\
T2 Half-Charge/E8 & Ladungs- und Adjungiertentypen genauer getrennt &
Renormierter quellenseitiger Half-Charge-Adapter mit Energie-, E8- und
Clockkontrolle \\
T3 gemeinsamer 3+1D-Ursprung & Komposition auf höheren Antworten
unterscheidbar & Ausgewählte lokale Teile, Nachbarschaft, Raumdimension
und gemeinsamer Parent \\
T4 chirale Materie & Minimaler Lorentz-Auslesekanal; Vertex-Händigkeit;
freier Nulltest & Nichtverschwindender nativer Feldadapter, chirales
Maß, Anomalien und Spiegelentkopplung \\
T5 Kontinuum/Dynamik & Wachsende Vergleichsfamilien und exakte
sektorübergreifende Gramformel & Größenuniforme physische Dynamik und
relativistischer wechselwirkender Limes \\
T6 Parameter/Spektren & Höhere Momente identifizieren zusätzliche
Kopplungen & Ursprung der Kopplungen und gemeinsame Zuordnung zu
gemessenen Massen und Konstanten \\
T7 Gravitation & Kein neuer Schließungsnachweis & Positiver dynamischer
Spin 2 und universelle Kopplung aus demselben Parent \\
T8 Zustand/Instrumente & Kohärenter Recorder, echter Eingriff und
Arbeitsbilanz & Ursprüngliche Präparation, autonome Kontrolle, Zeiger
und wiederverwendbare Aufzeichnung \\
\end{longtable}
}

Die Tabellen sind keine Vollständigkeitsbehauptung über sämtliche
offenen mathematischen Detailfragen. Sie lokalisieren den neuen Beitrag
innerhalb der erhaltenen Gesamtkarte. RH, Faktorisierung und P/NP wurden
in dieser Runde nicht untersucht; keine Aussage dieser Ausgabe schließt
sie.

\subsection{9. Der effektivste nächste
Versuch}\label{der-effektivste-nuxe4chste-versuch}

\subsubsection{A. Quellstrahlen zu gemeinsamem Träger - erster fehlender
Pfeil}\label{a.-quellstrahlen-zu-gemeinsamem-truxe4ger---erster-fehlender-pfeil}

Aus den tatsächlichen Originalobjekten sollen zwei Einbettungen
E\_L,E\_R sowie Bosonbank- und Vertexzuordnung folgen. Nicht nur eine
bekannte Überlappungszahl einsetzen, sondern Typen, Gruppenwirkung, CAR
und Operationen mitliefern. Das unmittelbar entscheidbare Ziel ist
\(E_L^\dagger E_R\), bei skalarem Ansatz cI.

\textbf{Akzeptanz:} Die Quelle wählt die Überlappung ohne Fit, und die
daraus berechnete globale Antwort besteht die I4-/Gramtests.
\textbf{Abbruch dieser behaupteten Herleitung:} Wenn mindestens zwei
unäquivalente Einbettungen unter denselben Quellbedingungen verbleiben,
wird die fehlende Auswahl als zusätzliche Annahme sichtbar belassen. Die
nächstkleinere zulässige Klassenbeschreibung wäre dann ein Ergebnis,
keine TOE-Schließung.

\subsubsection{B. Autonomer Eingriff statt gewähltem
Puls}\label{b.-autonomer-eingriff-statt-gewuxe4hltem-puls}

Den benötigten Nb-Eingriff und den Zeigerkoppler als Teile eines
gemeinsamen zeitunabhängigen Controller-System-Modells konstruieren. Es
muss die schon gerechnete Wirkung reproduzieren, Gesamtladung und
Gesamtenergie mitführen und den Zustand des Controllers vor und nach dem
Versuch angeben. Eine beliebig angesetzte Taktsequenz wäre noch keine
ursprüngliche Uhr.

\textbf{Akzeptanz:} Derselbe Anfangsvertrag berechnet beide
Interventionsarme, alle Zeigerzweige, Arbeit und Folgeantwort.
\textbf{Frühe Gegenprüfung:} Reine innere Symmetrieimpulse bei
kommutierendem Empfänger liefern nach dem bewiesenen Satz keinen solchen
Effekt.

\subsubsection{C. Einen nichtnullen, vertexkonsistenten Feldkanal auf
derselben
Quelle}\label{c.-einen-nichtnullen-vertexkonsistenten-feldkanal-auf-derselben-quelle}

Nicht den freien masselosen Nullansatz weiter ausbauen. Zuerst die
tatsächlichen Feldkomponenten zum Vertex und zur geladenen nativen
Antwort festlegen; dann das konkrete Matrixelement der Quelle unter der
vorgeschlagenen wechselwirkenden oder massiven Kinetik berechnen.

\textbf{Akzeptanz:} Nichtnuller Auslesekanal mit korrekter Ladung,
Adjunktion und Lorentztransformation; anschließend erst Pole,
Positivität und Skalierung. Ein algebraischer Projektor ohne nichtnulles
Matrixelement erfüllt das Ziel nicht.

Diese drei Wege können parallel bearbeitet werden, ihre Resultate müssen
anschließend denselben Träger und dieselben Ressourcen verwenden. Eine
große Raumzeitsimulation vor dieser Identifikation würde zusätzliche
Geometrieannahmen in das Ergebnis einbauen.

\subsection{10. Prüf- und
Veröffentlichungsstand}\label{pruxfcf--und-veruxf6ffentlichungsstand}

Vier neue Stränge wurden aus eingefrorenen Eingaben normal und optimiert
ausgeführt, mit Warnungen als Fehler und identischen Ergebnisbytes:
\textbf{514 eigene exakt arbeitende Bedingungen und vier numerische
Kontrollen je Variante}. Die exakte Kategorie enthält auch Quellenpins
und Strukturguards, nicht nur mathematische Sätze. Separat wurden 17
native Konstruktionsguards und 1073 alte P0/P1-Quellguards wiederholt.
Die gesamte v1.6.8-Suite wurde nicht erneut ausgeführt und wird nicht
als neue Verifikation gezählt.

Der Elternbeweis der komplexen Gramformel wurde zusätzlich unabhängig
gegen einen kleinen Bitmasken-CAR-Aufbau geprüft. Sein Modellumfang ist
eine Konventionsquerkontrolle, nicht ein zweiter vollständiger Lauf auf
dem großen nativen Träger. Die Voraussetzungen Frobeniusnorm 1,
Δ\textgreater0 und matrixaufgelöste Antwort werden ausdrücklich
beibehalten.

Alle neuen Prüfer, vollständigen Teilberichte, Quellenpins und die
bisherige vollständige Hauptfassung sind im versionierten Prüfpaket
enthalten. Hauptdokument, kurzes Update und einfache Fassung werden
zusammen in Documents abgelegt, ohne eine ältere Fassung zu
überschreiben. Dies ist keine Veröffentlichung, kein Commit/Push und
keine Änderung physischer Abschlussmarker.

\subsubsection{Automatisch aus dem abgeschlossenen
Replay}\label{automatisch-aus-dem-abgeschlossenen-replay}

{\def\LTcaptype{none} % do not increment counter
\begin{longtable}[]{@{}lrr@{}}
\toprule\noalign{}
Eigener Strang & Exakte Bedingungen & Numerische Bedingungen \\
\midrule\noalign{}
\endhead
\bottomrule\noalign{}
\endlastfoot
composition & 67 & 0 \\
field & 85 & 0 \\
protocol & 92 & 4 \\
source & 270 & 0 \\
\end{longtable}
}

Separat: 17 native Konstruktorguards, 1073 P0/P1-Quellguards und 13
unabhängige kleine CAR-Querkontrollen. Quellenpins und Strukturguards
sind in der exakten Kategorie enthalten. Die gesamte v1.6.8-Suite wurde
nicht neu ausgeführt.

\begin{center}\rule{0.5\linewidth}{0.5pt}\end{center}

\section{Vollständiger Beweis: höhere Antwort und positive
Kopplungsmatrix}\label{vollstuxe4ndiger-beweis-huxf6here-antwort-und-positive-kopplungsmatrix}

\section{Eine höhere Antwort liest die Zusammensetzung
aus}\label{eine-huxf6here-antwort-liest-die-zusammensetzung-aus}

\begin{enumerate}
\def\labelenumi{\arabic{enumi}.}
\setcounter{enumi}{14}
\tightlist
\item
  September 2026 · v1.6.9 · vollständiger eigener Kompositionsbeweis
\end{enumerate}

\subsection{1. Frage und Ergebnis}\label{frage-und-ergebnis}

Kann die innere W-Struktur schon bestimmen, wie mehrere Vorkommen
gemeinsam wirken? Für eine ausdrücklich erklärte, besonders einfache
Familie lässt sich die Frage exakt beantworten: \textbf{Nein, nicht aus
W und auch nicht aus dem gesamten Zweiladungsspektrum allein.} Die
Dreiladungsantwort bestimmt jedoch genau die zusätzlich gewählte
positive Kopplungsmatrix. Das ist ein konstruktiver Anschluss für die
vom Benutzer vorgeschlagene Rückrechnung aus „Schatten``.

Es handelt sich um zwei Aussagen innerhalb einer festgelegten
Modellklasse, nicht um eine allgemeine Unmöglichkeit jeder
TFPT-Herleitung. Der globale Kopplungstensor ist hier
Gegenmodelleingabe; seine compilerseitige Auswahl wurde nicht gefunden.
Beide Gegenmodelle haben einen gemeinsamen globalen Träger und echte
globale Hamiltonoperatoren.

\subsection{2. Gemeinsame ursprüngliche
Daten}\label{gemeinsame-urspruxfcngliche-daten}

Wir verwenden den unveränderten realen Tensor W mit 60 Kanälen und 2016
Fermionpaarspalten. Aus jeder Zeile wird eine antisymmetrische
64-mal-64-Matrix A\_A: oberhalb der Diagonale steht W, unterhalb sein
negatives Transponiertes. Direkt aus allen 480 Vertices folgt

\[\adjustbox{max width=\linewidth}{$\displaystyle 
\tfrac12\operatorname{tr}(A_A A_B^\dagger)=8\delta_{AB},\qquad
\sum_A A_A^\dagger A_A=15I_{64}.
$}\]

Die native Dreifermion-zu-Boson-Fermion-Abbildung heiße \(T_0\). Ihr
Gramoperator ist \(G_0=T_0T_0^\dagger\) auf 3840 Dimensionen. Sämtliche
geordneten Fermiontripel und CAR-Vorzeichen werden neu aufgebaut. Es
ergeben sich 29760 nichtverschwindende Einträge und das exakt
kontrollierte Spektrum

\[\adjustbox{max width=\linewidth}{$\displaystyle 
\operatorname{spec}(G_0)=\{0^{[64]},7^{[2880]},10^{[576]},12^{[320]}\}.
$}\]

Die Hochzahlen in eckigen Klammern sind Multiplizitäten. Der Prüfer
kontrolliert das ganzzahlige annihilierende Polynom
\(G_0(G_0-7I)(G_0-10I)(G_0-12I)=0\) und die rationalen
Spektralprojektorspuren. Eine Zeilensummen-Normschranke sichert die
benutzten int64-Produkte und Spuren gegen Überlauf ab; keine
Gleitkomma-Eigenwerte begründen diese Aussage.

Setze

\[\adjustbox{max width=\linewidth}{$\displaystyle 
K=8I_{3840}-G_0,\qquad \operatorname{tr}K=960.
$}\]

Unabhängig wird für alle 3600 Kanalblöcke exakt geprüft:

\[\adjustbox{max width=\linewidth}{$\displaystyle 
K_{Ai,Bj}=(A_B A_A^\dagger)_{ij}.
$}\]

\subsection{3. Ein zusätzlicher symmetrischer
Kompositionstensor}\label{ein-zusuxe4tzlicher-symmetrischer-kompositionstensor}

Fermionen erhalten einen zusätzlichen L-dimensionalen
Multiplizitätsfaktor; die 60 Bosonen bleiben \textbf{eine einzige
gemeinsame Bank}. Wähle eine komplex-symmetrische Matrix C mit
\(\operatorname{tr}(C^\dagger C)=1\). Die antisymmetrischen globalen
Paarmatrizen lauten

\[\adjustbox{max width=\linewidth}{$\displaystyle 
M_A=C\otimes A_A,\qquad
P_A(C)=\sum_{(x,i)<(y,j)}(M_A)_{xi,yj}f_{yj}f_{xi},
$}\] \[\adjustbox{max width=\linewidth}{$\displaystyle 
H_C=\Delta N_b+g\sum_A\bigl(b_A^\dagger P_A(C)+P_A(C)^\dagger b_A\bigr).
$}\]

Die innere Gruppe wirkt auf dem 64-Faktor, trivial auf der
Multiplizität. Daher bleiben die nativen inneren Intertwiner, der
diagonal geliftete Clock und \(N=N_f+2N_b\) erhalten. Die vollständige
globale Paarabbildung hat für \textbf{jedes} solche C den Gramoperator
\(8I_{60}\).

Wichtig: Hier bleibt die global normierte Paarantwort gleich. Die
einzelne Kopie trägt im Allgemeinen \textbf{nicht} jeden lokalen
W-Koeffizienten mit unveränderter Stärke g. Genau darin unterscheidet
sich diese Familie vom Zwei-Chart-Modell des Quellenstrangs, das jede
lokale W-Kopplung unverändert lässt, dafür aber zwei Bosonbanken hat.
Beide Klassen prüfen unterschiedliche Eindeutigkeitsbehauptungen und
werden nicht vermischt.

\subsection{4. Gleiche vollständige Zweiladungsspektren,
unterschiedliche globale
Dynamik}\label{gleiche-vollstuxe4ndige-zweiladungsspektren-unterschiedliche-globale-dynamik}

Zwei einfache Familien sind

\[\adjustbox{max width=\linewidth}{$\displaystyle 
C_{\rm verteilt}=\frac{I_L}{\sqrt L},\qquad
C_{\rm kollektiv}=\frac{J_L}{L},
$}\]

wobei J\_L die Matrix aus lauter Einsen bezeichnet. Beide sind
symmetrisch, normiert, invariant unter Permutation der Kopien und
stimmen für L=1 mit dem nativen Modell überein. Sie liefern jedoch
verschiedene spätere Physik. Permutationssymmetrie ist dabei keine
geometrische Lokalität.

Auf N=2 besteht H\_C aus 60 identischen hellen Zweierblöcken mit
Kopplung \(\sqrt8g\), ergänzt um \(\binom{64L}{2}-60\) dunkle
Paarzustände. Bei festem L sind deshalb die \textbf{gesamten}
N=2-Hamiltonspektren einschließlich sämtlicher Multiplizitäten gleich,
nicht nur die hellen Pole. Die Bezeichnung „gleiches Spektrum`` enthält
keine Gleichheit aller markierten Präparationen und Instrumente.

Für L=2 hat dieser Sektor 8188 Dimensionen: 8128 Fermionpaare und 60
Bosonen, mit 8068 dunklen Paarzuständen. Die primitive Kopplungsmetrik
ist dagegen

\[\adjustbox{max width=\linewidth}{$\displaystyle 
\sum_A M_A^\dagger M_A=15(C^\dagger C)\otimes I_{64}.
$}\]

Der verteilte Fall hat keinen vollständig abgekoppelten primitiven
Fermionmodus, der kollektive Fall hat 64 davon. Für allgemeines L sind
es im kollektiven Fall 64(L−1). Dies zeigt bereits eine andere
CAR-Kopplungsstruktur. Ein stärkerer, nicht auf markierte Einzelmoden
angewiesener Zeuge ist das folgende vollständige N=3-Spektrum.

